The implementation is organized as a FastAPI backend and React frontend. The backend exposes route ranking through the same workflow used by the evaluation scripts: graph construction, route generation, feature extraction, profile construction, ranking, and result serialization. The OSM pseudo-history files, quality reports, baseline outputs, and Porto benchmark outputs are written under the project data directory. The public Porto CSV is intentionally not committed because of its size; the evaluator writes a status report if the dataset is missing.

The main OSM route-candidate evaluation can be reproduced with the route-candidate baseline script. The public Porto benchmark can be reproduced after placing the Porto training CSV at \texttt{data/porto/train.csv}. The 100-query run reported in this paper uses:

\begin{verbatim}
python scripts\evaluate_porto_candidate_baselines.py
  --max-tests 100 --max-rows-to-scan 30000
  --sample-stride 20 --k-routes 10
  --dist-meters 2500 --bootstrap-samples 1000
  --min-anchor-spacing-m 120 --max-anchors 50
  --shared-graph-radius-m 8000
\end{verbatim}

The script outputs both JSON and CSV summaries, including skipped-baseline reasons. Prompt/SBERT and hybrid are skipped for Porto because the public trajectories do not contain natural-language preferences. This behavior is deliberate: missing data should produce a reported skip, not fabricated preference labels.
